\documentclass[10pt,a4paper]{article}

\usepackage[utf8]{inputenc}
\usepackage{graphicx}
\usepackage{cite}
\usepackage{hyperref}
\usepackage{geometry}
\usepackage{changepage}
\usepackage{booktabs}
\usepackage{amsmath}
\usepackage{fancyhdr}
\usepackage{graphicx}
\usepackage{longtable}
\usepackage{array}

\pagestyle{fancy}
\fancyhf{}
\fancyfoot[C]{\thepage}

\title{Related Works on Personalized Healthcare using Artificial Intelligence}
\author{Djelloul Daouadji Fadela}

\begin{document}

\maketitle

\begin{abstract}
This report summarizes recent research and developments in the field of AI-based personalized healthcare. It highlights the key contributions, methodologies, and challenges faced by existing systems.
\end{abstract}

\section{Introduction}
The use of Artificial Intelligence (AI) in healthcare has opened new avenues for personalized treatment. This report reviews existing studies and methods used in AI-driven healthcare solutions.

\section{Related Works}
\input{./related_work/artificial_intelligence}
\input{./related_work/reinforcement_learning}

% \input{./related_work/deep_learning}
% 

% ============================== COMPARISON TABLE ==============================
\section{Comparison of the Solutions}
\begin{table}[htbp]
\begin{adjustwidth}{-2cm}{-2cm}
\centering
\setlength{\tabcolsep}{4pt}
\renewcommand{\arraystretch}{0.8} 
\scriptsize
\begin{tabular}{|p{1.5cm}|p{2.5cm}|p{3cm}|p{3.5cm}|p{2.5cm}|p{2.5cm}|}
\hline
\textbf{Work} & \textbf{Disease/Domain} & \textbf{Dataset} & \textbf{Data Processing} & \textbf{Approach} & \textbf{Results} \\
\hline
\cite{dejongRealizingVisionPrecision2021} & Epilepsy & Phase III (235) + Validation (47) patients & Not specified & Gradient Boosting Decision Tree (GBDT) & AUC: 0.76 (train), 0.75 (validation) \\
\hline
\cite{tasinDiabetesPredictionUsing2023} & Diabetes Prediction & Pima Indian (768) + RTML (203) records & Imputation, ADASYN, Mutual Info, Holdout Validation & XGBoost + Ensemble Methods (voting, bagging) & AUC: 0.84, Accuracy: 81\%, F1 Score: 0.81 \\
\hline
\cite{sarkarIntegratingMachineLearning2025} & Alzheimer's Disease & Wearable sensors and motion capture data & Normalization, median imputation, SMOTE, RFE, correlation analysis & Hybrid CNN-RNN (LSTM) & Accuracy: 93\%, Precision: 92\%, Recall: 91\%, F1-Score: 91.5\%, AUC-ROC: 95\% \\
\hline
\cite{swapnaDiabetesDetectionUsing2018} & Diabetes & ECG recordings (71 datasets) & Pan-Tompkins for QRS detection & CNN-LSTM + SVM & Accuracy: 95.7\% \\
\hline
\cite{shuvoCardioXNetNovelLightweight2021} & Cardiovascular Disease & PCG datasets (GitHub + PhysioNet) & Normalization & Lightweight CRNN (CardioXNet) & Accuracy: 99.6\% (GitHub), 86.57\% (PhysioNet) \\
\hline
\cite{anjumOptimizingType22024} & Type 2 Diabetes & CGM data from 8 patients & NaN removal, feature extraction, dimensionality reduction, 75/25 split & XGBoost, SARIMA, Prophet regression for BG prediction & Not specified \\
\hline
\cite{wangDeepLearningAlgorithm2021} & COVID-19 diagnosis & 1,065 CT images & ROI extraction, 299x299 resize, RGB conversion & M-Inception, fine-tuned transfer learning & 89.5\% internal validation; 85.2\% on difficult cases \\
\hline
\cite{kermanyIdentifyingMedicalDiagnoses2018} & Medical imaging diagnostics & OCT (108,312) + X-rays (5,856) & Not specified & Transfer learning with pre-trained CNN & 96.6\% (OCT), 92.8\% (pneumonia) \\
\hline
\cite{gayathriEnhancingHeartDisease2024} & Heart Disease & UCI Cleveland Heart Disease dataset & Data augmentation & Reinforcement Learning (RL) & Accuracy: 94\% \\
\hline
\cite{javadReinforcementLearningBased2019} & Type 1 Diabetes & Clinical data from 87 patients (MGH, 2003--2013) & Discretization & Q-Learning (model-free RL) & Accuracy: 88\% \\
\hline
\cite{abdelazizDeepQNetworkDQN2025} & Heart Disease & UCI Heart Disease Dataset via Kaggle & Stratified 80/20 split, robust scaling & Deep Q-Network (DQN) & Accuracy: 98.41\%, MSE: 0.0001 \\
\hline
\cite{barataReinforcementLearningModel2023} & Skin Cancer & HAM10000 dataset (10,000+ dermatoscopic images) & Feature extraction via CNN, normalization & RL (Deep Q-learning), SL (CNN) & Not specified \\
\hline
\cite{stemberReinforcementLearningUsing2022} & Brain Tumor Localization & BraTS 2014 (T1-weighted MRI, 60 images) & Grid-based image division, 60x60 pixel navigation & Deep Q-Network (DQN) & 70\% accuracy vs 11\% for SL \\
\hline
\cite{petersenDeepReinforcementLearning2019} & Sepsis Treatment & 500 simulated patients using IIRABM & Not specified & Deep Reinforcement Learning (DRL) & Not specified \\
\hline
\cite{yangPrescDRLDeepReinforcement2024} & Diabetes Treatment & Custom benchmark dataset & Feature extraction & Deep Reinforcement Learning (DRL) & 117-153\% improvement in rewards, 40.5\% precision, 63\% recall \\
\hline
\cite{yangDeepReinforcementLearning2024} & Multi-class Clinical Classification & Real-world clinical datasets & Feature selection, normalization & Combined Dueling and Double Deep Q-learning & Not specified \\
\hline
\cite{maDeepAttentionQnetwork2023} & Sepsis and Acute Hypotension & Real-world datasets & Not specified & Transformer-based Deep Attention Q-Network & Better expected rewards than clinician policies \\
\hline
\cite{tardiniOptimalTreatmentSelection2022} & Oropharyngeal squamous cell carcinoma & 536 patient records (MD Anderson, 2005--2013) & 75/25 split; radiomics in select training runs & Deep Q-Learning & 87.35\% accuracy; +3.73\% survival, -0.75\% dysphagia \\
\hline
\cite{zhengPersonalizedMultimorbidityManagement2021} & Type 2 Diabetes & 16,665 EHRs from NYU Langone (2009--2017) & 60/40 split; excluded non-longitudinal visits & Reinforcement Learning (RL) models & High clinician concordance; better outcomes (86.1\% glycemia, 98.4\% CVD) \\
\hline
\cite{liuDeepReinforcementLearning2022} & Cancer therapy personalization & 1080 RNA-seq patients, 223 drugs, TCGA + GDSC & Filtered by PubChem ID & PPORank (DRL), Actor-Critic with PPO & NDCG close to 0.8 \\
\hline
\cite{zhaoSafeenhancedFullyClosedloop2025} & Type 1 Diabetes (Artificial Pancreas) & Simglucose simulator (30 virtual patients) & State definition, reward engineering, normalization & Enhanced PPO (DRL) with dual safety & Median TIR: 87.45\%; eliminated severe hypoglycemia \\
\hline
\cite{apostolopoulosCovid19AutomaticDetection2020} & COVID-19 Detection (X-Ray) & Combined X-ray datasets (COVID-19, other pneumonia, healthy) & Resizing, normalization, data augmentation, labeling & Transfer Learning with CNNs (e.g., VGG, ResNet) & High accuracy (typically $>$90\%); rapid screening \\
\hline
\cite{liuReinforcementLearningOptimize2024} & Mechanical Ventilation Optimization & Retrospective ICU patient EHR data (Physiological/Ventilator settings) & Cleaning, imputation, feature engineering, state/action/reward definition & Offline Reinforcement Learning (RL) & Improved physiological parameters; potential for reduced ventilation duration \\
\hline

\end{tabular}
\caption{Comparison of AI Approaches in Health Applications}
\label{tab:ai_health_comparison}
\end{adjustwidth}
\end{table}

\section{Conclusion}
Personalized healthcare using AI continues to evolve, offering significant potential to improve patient care. However, integration into real-world clinical settings remains an ongoing challenge.

\newpage
% Bibliography
\bibliographystyle{unsrt}
\bibliography{bibliography}

\end{document}
